% GENERATED FILE. Edit canonical lessons, then run npm run generate:books.
% canonical-source-hash: fnv1a64:439f13635821241c
% canonical-lessons: TE-C450-parkingu, TE-C450-nilupu, TE-C450-ekkincu, TE-C450-dincu, TE-C450-digabettu

\chapter{Parking and Dropping Off}
\label{ch:te-parking-and-dropping-off}
\hlchaptermodality{450}

\begin{chapteropening}
\textbf{By the end of this chapter:} \emph{I can say parking, to pull up, to put aboard, to drop off and to drop someone off.}

Complete the last lesson of chapter 450: I can say parking, to pull up, to put aboard, to drop off and to drop someone off.
\end{chapteropening}

\section[\texorpdfstring{\te{పార్కింగు} (pārkiṅgu)}{పార్కింగు (pārkiṅgu)}]{\te{పార్కింగు} (pārkiṅgu) — parking}

\label{lesson:TE-C450-parkingu}



\begin{quote}
Before the new one: say the Telugu for to seize up, then the Telugu for to get blocked.
\end{quote}

\subsection*{You'll want to know: \te{పార్కింగు}}
\textbf{\te{పార్కింగు}} — \emph{pārkiṅgu} — \textquotedblleft{}parking\textquotedblright{}.

\begin{quote}
\te{ఆ} \te{దుకాణం} \te{దగ్గర} \te{పార్కింగు} \te{లేదు}.
\end{quote}

\begin{quote}
\emph{ā dukāṇaṁ daggara pārkiṅgu lēdu.}
\end{quote}

\begin{quote}
There is no parking near that shop.
\end{quote}

A parking place is \te{పార్కింగు} \te{స్థలం}.

\begin{cousinweb}
\textbf{\te{పార్కింగు}} (\emph{pārkiṅgu}) is a loan from English \emph{parking}, with the Telugu final -\te{ఉ}.
\end{cousinweb}

\begin{grammarlens}[title={What you've built}]
One more word for getting around.
\end{grammarlens}

\subsection*{Your turn}
\begin{itemize}
\raggedright
  \item \emph{Say it:} \emph{pārkiṅgu}
  \item \emph{Say it:} \emph{pārkiṅgu}, once more
  \item \emph{Say it:} \emph{ā dukāṇaṁ daggara pārkiṅgu lēdu.} — There is no parking near that shop.
  \item \emph{Recall:} say the Telugu for to seize up, then the Telugu for to get blocked, then say \emph{pārkiṅgu} again
\end{itemize}

\begin{tcolorbox}[breakable,colback=teal!4,colframe=teal!35!black,title={Before you move on}]
What does \te{పార్కింగు} mean? (\textquotedblleft{}Parking\textquotedblright{}.) Say it once more.
\end{tcolorbox}
\section[\texorpdfstring{\te{నిలుపు} (nilupu)}{నిలుపు (nilupu)}]{\te{నిలుపు} (nilupu) — to pull up; to park}

\label{lesson:TE-C450-nilupu}



\begin{quote}
Before the new one: say the Telugu for to get blocked, then the Telugu for parking.
\end{quote}

\subsection*{You'll want to know: \te{నిలుపు}}
\textbf{\te{నిలుపు}} — \emph{nilupu} — \textquotedblleft{}to pull up; to park\textquotedblright{}.

\begin{quote}
\te{అతను} \te{కారును} \te{చెట్టు} \te{కింద} \te{నిలిపాడు}.
\end{quote}

\begin{quote}
\emph{atanu kārunu ceṭṭu kinda nilipāḍu.}
\end{quote}

\begin{quote}
He parked the car under a tree.
\end{quote}

Past (I …): \te{నిలిపాను}, \emph{nilipānu}.

The causative of \te{నిలుచు}, to stand: to make something stand. It also means to stop a process.

\begin{cousinweb}
\textbf{\te{నిలుపు}} (\emph{nilupu}) is native Dravidian: DEDR 3675 lists Telugu \emph{nilupu}, \textquotedblleft{}to cause to stand, to stop\textquotedblright{}, beside \emph{nilucu}, \textquotedblleft{}to stand\textquotedblright{}.
\end{cousinweb}

\begin{grammarlens}[title={What you've built}]
One more word for getting around.
\end{grammarlens}

\subsection*{Your turn}
\begin{itemize}
\raggedright
  \item \emph{Say it:} \emph{nilupu}
  \item \emph{Say it:} \emph{nilupu}, once more
  \item \emph{Say it:} \emph{atanu kārunu ceṭṭu kinda nilipāḍu.} — He parked the car under a tree.
  \item \emph{Recall:} say the Telugu for to get blocked, then the Telugu for parking, then say \emph{nilupu} again
\end{itemize}

\begin{tcolorbox}[breakable,colback=teal!4,colframe=teal!35!black,title={Before you move on}]
What does \te{నిలుపు} mean? (\textquotedblleft{}To pull up; to park\textquotedblright{}.) Say it once more.
\end{tcolorbox}
\section[\texorpdfstring{\te{ఎక్కించు} (ekkiñcu)}{ఎక్కించు (ekkiñcu)}]{\te{ఎక్కించు} (ekkiñcu) — to put aboard; to load}

\label{lesson:TE-C450-ekkincu}



\begin{quote}
Before the new one: say the Telugu for parking, then the Telugu for to pull up.
\end{quote}

\subsection*{You'll want to know: \te{ఎక్కించు}}
\textbf{\te{ఎక్కించు}} — \emph{ekkiñcu} — \textquotedblleft{}to put aboard; to load\textquotedblright{}.

\begin{quote}
\te{నాన్న} \te{పిల్లలను} \te{బస్సు} \te{ఎక్కించాడు}.
\end{quote}

\begin{quote}
\emph{nānna pillalanu bassu ekkiñcāḍu.}
\end{quote}

\begin{quote}
Father put the children on the bus.
\end{quote}

Past (I …): \te{ఎక్కించాను}, \emph{ekkiñcānu}.

The causative of \te{ఎక్కు}, to climb on, to board: to make someone or something get on.

\begin{cousinweb}
\textbf{\te{ఎక్కించు}} (\emph{ekkiñcu}) is native Dravidian: DEDR 766 lists Telugu \emph{ekkincu}, \textquotedblleft{}to cause to ascend\textquotedblright{}, beside \emph{ekku}, \textquotedblleft{}to climb, to get on\textquotedblright{}.
\end{cousinweb}

\begin{grammarlens}[title={What you've built}]
One more word for getting around.
\end{grammarlens}

\subsection*{Your turn}
\begin{itemize}
\raggedright
  \item \emph{Say it:} \emph{ekkiñcu}
  \item \emph{Say it:} \emph{ekkiñcu}, once more
  \item \emph{Say it:} \emph{nānna pillalanu bassu ekkiñcāḍu.} — Father put the children on the bus.
  \item \emph{Recall:} say the Telugu for parking, then the Telugu for to pull up, then say \emph{ekkiñcu} again
\end{itemize}

\begin{tcolorbox}[breakable,colback=teal!4,colframe=teal!35!black,title={Before you move on}]
What does \te{ఎక్కించు} mean? (\textquotedblleft{}To put aboard; to load\textquotedblright{}.) Say it once more.
\end{tcolorbox}
\section[\texorpdfstring{\te{దించు} (diñcu)}{దించు (diñcu)}]{\te{దించు} (diñcu) — to unload; to lower}

\label{lesson:TE-C450-dincu}



\begin{quote}
Before the new one: say the Telugu for to pull up, then the Telugu for to put aboard.
\end{quote}

\subsection*{You'll want to know: \te{దించు}}
\textbf{\te{దించు}} — \emph{diñcu} — \textquotedblleft{}to unload; to lower\textquotedblright{}.

\begin{quote}
\te{కూలీలు} \te{లారీ} \te{నుంచి} \te{బస్తాలు} \te{దించారు}.
\end{quote}

\begin{quote}
\emph{kūlīlu lārī nuñci bastālu diñcāru.}
\end{quote}

\begin{quote}
The labourers unloaded the sacks from the lorry.
\end{quote}

Past (I …): \te{దించాను}, \emph{diñcānu}.

The causative of \te{దిగు}, to get down: to bring something down. \te{బరువు} \te{దించు}, to put down a load.

\begin{cousinweb}
\textbf{\te{దించు}} (\emph{diñcu}) is native Dravidian: DEDR 502 lists Telugu \emph{dincu} beside \emph{digu}, \textquotedblleft{}to descend, to get off\textquotedblright{}.
\end{cousinweb}

\begin{grammarlens}[title={What you've built}]
One more word for getting around.
\end{grammarlens}

\subsection*{Your turn}
\begin{itemize}
\raggedright
  \item \emph{Say it:} \emph{diñcu}
  \item \emph{Say it:} \emph{diñcu}, once more
  \item \emph{Say it:} \emph{kūlīlu lārī nuñci bastālu diñcāru.} — The labourers unloaded the sacks from the lorry.
  \item \emph{Recall:} say the Telugu for to pull up, then the Telugu for to put aboard, then say \emph{diñcu} again
\end{itemize}

\begin{tcolorbox}[breakable,colback=teal!4,colframe=teal!35!black,title={Before you move on}]
What does \te{దించు} mean? (\textquotedblleft{}To unload; to lower\textquotedblright{}.) Say it once more.
\end{tcolorbox}
\section[\texorpdfstring{\te{దిగబెట్టు} (digabeṭṭu)}{దిగబెట్టు (digabeṭṭu)}]{\te{దిగబెట్టు} (digabeṭṭu) — to drop someone off}

\label{lesson:TE-C450-digabettu}



\begin{quote}
Before the new one: say the Telugu for to put aboard, then the Telugu for to unload.
\end{quote}

\subsection*{You'll want to know: \te{దిగబెట్టు}}
\textbf{\te{దిగబెట్టు}} — \emph{digabeṭṭu} — \textquotedblleft{}to drop someone off\textquotedblright{}.

\begin{quote}
\te{నా} \te{స్నేహితుడు} \te{నన్ను} \te{స్టేషన్} \te{దగ్గర} \te{దిగబెట్టాడు}.
\end{quote}

\begin{quote}
\emph{nā snēhituḍu nannu sṭēṣan daggara digabeṭṭāḍu.}
\end{quote}

\begin{quote}
My friend dropped me off near the station.
\end{quote}

Past (I …): \te{దిగబెట్టాను}, \emph{digabeṭṭānu}.

\te{దిగు}, to get down, + \te{పెట్టు}, to put: to see someone safely to a place.

\begin{cousinweb}
\textbf{\te{దిగబెట్టు}} (\emph{digabeṭṭu}) is native Dravidian: \te{దిగు}, \textquotedblleft{}to get down\textquotedblright{} (DEDR 502), and \te{పెట్టు}, \textquotedblleft{}to put, to place\textquotedblright{} (DEDR 4390).
\end{cousinweb}

\begin{grammarlens}[title={What you've built}]
One more word for getting around.
\end{grammarlens}

\subsection*{Your turn}
\begin{itemize}
\raggedright
  \item \emph{Say it:} \emph{digabeṭṭu}
  \item \emph{Say it:} \emph{digabeṭṭu}, once more
  \item \emph{Say it:} \emph{nā snēhituḍu nannu sṭēṣan daggara digabeṭṭāḍu.} — My friend dropped me off near the station.
  \item \emph{Recall:} say the Telugu for to put aboard, then the Telugu for to unload, then say \emph{digabeṭṭu} again
\end{itemize}

\begin{tcolorbox}[breakable,colback=teal!4,colframe=teal!35!black,title={Before you move on}]
What does \te{దిగబెట్టు} mean? (\textquotedblleft{}To drop someone off\textquotedblright{}.) Say it once more.
\end{tcolorbox}
